\section{Audit findings, status changes, and reproducibility}
\label{sec:audit}

\subsection{What has changed mathematically}

The source snapshot matters because several incoming reports overlap,
and some questions in an earlier report are answered by a later one.
We inspected the canonical manuscript and the textual contents of all
six incoming archives at the pinned revision. We used that comparison
to choose the present targets and to identify the inherited inputs.
This is an audit of the dependencies and claims relevant to this
continuation, not a new line-by-line certification of all 379 pages.

The main changes of status are as follows.

\begingroup
\small
\renewcommand{\arraystretch}{1.18}
\begin{longtable}{@{}>{\raggedright\arraybackslash}p{.28\textwidth}
 >{\raggedright\arraybackslash}p{.31\textwidth}
 >{\raggedright\arraybackslash}p{.35\textwidth}@{}}
\toprule
Question in the snapshot & Previous status & Result in this article\\
\midrule
\endfirsthead
\toprule
Question in the snapshot & Previous status & Result in this article\\
\midrule
\endhead
Global limit of the optimal kernel constants
& Local scaling theorem; global limit conjectured in
  \cite[Section 8]{IncomingSharp}
& Theorem~\ref{ker:thm:limit} proves the limit, the limiting maximizer,
  and a uniform quantitative enclosure.\\
Uniqueness and decrease of kernel maxima
& Conjectured at every finite truncation
& Theorem~\ref{ker:thm:N2} proves uniqueness at $N=2$;
  Theorem~\ref{ker:thm:analytic} proves eventual uniqueness,
  decrease, and a convergent expansion.\\
Sharp decay of $C_N-1$
& A lower asymptotic; matching upper asymptotic conjectured in
  \cite[Section 9]{IncomingBounds}
& Theorems~\ref{ord:thm:axis-reduction} and
  \ref{ord:thm:leading} prove the global asymptotic and the eventual
  axis reduction.\\
Next order scale and near maximizers
& Not resolved by the leading candidate
& Theorem~\ref{ord:thm:second} gives the threshold-layer correction;
  Corollary~\ref{ord:cor:near-maximizers} gives the exact
  $o(N^{-p})$ criterion.\\
Polynomial Appell root separation
& Exponential lower bound; a polynomial bound requested in
  \cite[Section 8]{IncomingUniform}
& Theorem~\ref{mesh:thm:polynomial} proves
  $\mesh(Q_n)\ge1/(10n^2)$.\\
Growing-index Lerch saturation
& Explicit exponential threshold and logarithmic index range
& Theorem~\ref{mesh:thm:saturation} gives an
  $O(n^4\log^2n)$ threshold and a uniform $8/\sqrt{k}$ log error.\\
Signs at outer orders $-2,-3$
& Individual opposite-sign evaluations; negativity proved down to $-1$
& Theorem~\ref{neg:thm:transitions} classifies all $b>0$:
  one simple crossing in each case.\\
Mixed Gaussian periods
& $S_4$ proved; $S_6,S_8,S_{10},S_{12}$ conjectural
& Exact generating and sampling identities are proved.
  A new $S_{14}$ candidate is supplied with certified proximity;
  its equality remains conjectural.\\
\bottomrule
\end{longtable}
\endgroup

The constants $M_N$ and $C_N$ must remain distinct. The former is the
maximum of a pointwise two-variable kernel; the latter is the supremum
after averaging that kernel against the particular gamma distributions
coming from the two polylogarithm orders. Although $C_N\le M_N$, their
limits are different:
\[
 M_N\longrightarrow1.05386193036\ldots,\qquad C_N\longrightarrow1.
\]
The limiting pointwise extremizer has $Np_N\to c_\infty$ and
$y_N\to y_\infty$. The order-averaged extremizer instead approaches
the outer-order boundary, with $b_N\sim\log N/\log(3/2)$. The probability
laws allowed by the order parameters do not reproduce an arbitrary
point mass at the pointwise maximizing pair.

\subsection{Proposed corrections and precise integration instructions}

No theorem-level error was established in the inherited results used
in these proofs. The necessary manuscript changes are chiefly updates
of status and explicit qualifications. They should not be presented
as retractions of correct earlier results.

\paragraph{The two Euler conjectures.}
In \cite{IncomingSharp}, replace the conjectural status of
\texttt{conj:limit} by Theorem~\ref{ker:thm:limit}. Retain the
all-finite-index content of \texttt{conj:maxima}: the present eventual
theorem and the exact $N=2$ case leave intermediate finite indices
unclassified. In \cite{IncomingBounds}, replace
\texttt{conj:euler-rate} by Theorem~\ref{ord:thm:leading}, and add
the exact axis-reduction theorem before discussing the rate.
It supplies the global upper bound that an axis computation alone
cannot give.

The strict domain-specific estimate $R_N(a,b)<1$ for $a\ge1$ in the
canonical \path{chapters/05-real-euler.tex} is an explicit dependency.
The larger universal constant $C_*$ does not supply the same
exclusion. The already proved value
$C_*=1.1365611033395\ldots$ remains a separate, all-index optimum.

\paragraph{Quantifiers on asymptotics.}
The uniqueness and monotonicity statements for the two maximizing
problems are eventual unless explicitly stated at $N=2$.
We do not provide a numerical starting index for those eventual
statements. Also, the next Taylor moment of the axis kernel is not
the next asymptotic term of the optimized remainder.
The threshold layer gives
$N^{-\kappa}(\log N)^{-1/2}$, and the first relative inverse-log
coefficient is large in magnitude. Remark~\ref{ord:rem:late} explains
why moderate-index agreement is an unreliable way to assess this
second term. This is a qualification needed in further extrapolation,
not a claim that the incoming reports made the incorrect expansion.

\paragraph{Lerch geometry.}
In \cite[Section 4]{IncomingUniform}, retain the original mesh bound:
it is useful at small degrees and is compatible with the stronger
polynomial estimate. Replace an exponential threshold as the only
available effective theorem by
Theorems~\ref{mesh:thm:saturation} and \ref{mesh:thm:transfer}.
The sharper generic transfer can use any proved mesh bound, including
the maximum of the old and new lower bounds.
The family at $\rho=1$ is the derivative family with $k\ge1$;
the undifferentiated Lerch sum need not converge there.
The polynomial growing-index theorem establishes leading relative
locations and the full zero count, not uniform constant offsets or
complete zero expansions in the joint limit.

\paragraph{Negative outer orders.}
Keep the previously proved negativity for every real outer order
at least $-1$. The two formulas at $b=1$ used in
Theorem~\ref{neg:thm:transitions} were already in the incoming Euler
report and should retain that provenance. What is new is the unique
simple transition throughout $b>0$. The kernel $q_5$ shows why the
one-crossing proof cannot simply be asserted for all lower orders.

\paragraph{Arithmetic status and old labels.}
The canonical \path{chapters/04-S4-proof.tex} labels
\texttt{s4proof:thm:s4} and \texttt{gaussian:thm:S4} denote an
already proved result; the older \texttt{gaussian:conj:S4} survives
as a compatibility label, not as evidence that the identity is still
open. By contrast, \texttt{cycloquot:conj:S6} in
\path{chapters/04-cyclotomic-quotients.tex} and
\texttt{s8new:conj:S8} in \path{chapters/04-S8-candidate.tex}
remain conjectures. The $S_{10},S_{12}$ vectors in
\cite[Section 5]{IncomingUniform} are prior candidates.
Our weight-fifteen vector is new relative to those files; restricted
relation-space calculations at the same weight were already present.

A rational enclosure containing zero proves a bound on the residual,
not its exact vanishing. Likewise, nonmembership in a specified
finite relation space does not prove independence of complex periods
or impossibility of a reduction using a larger alphabet.
The generating identities of Section~\ref{mix:sec} also do not
authorize coefficientwise fixed-weight conclusions beyond what their
proofs actually establish.

The companion integration note gives the source files and labels in
machine-searchable form. Labels in this article use the prefixes
\texttt{ker:}, \texttt{ord:}, \texttt{mesh:}, \texttt{neg:}, and
\texttt{mix:} to reduce collision risk when sections are moved.

\subsection{What the computational artifacts prove}

The package separates exact replay from numerical diagnostics.
The distinction is mathematical: increasing working precision is not
a substitute for bounding an infinite tail, and a computed optimizer
is not a global maximum certificate.

\begin{enumerate}
\item \textbf{Exact rational kernel enclosures.}
  \path{code/verify_kernel.py} isolates the root defining $y_\infty$
  using rational bounds for the logarithm, encloses the corresponding
  exponential constant, and isolates the algebraic $N=2$ root and
  its rationally defined maximum. The qualitative uniqueness and
  global conclusions are supplied by Section~\ref{sec:kernel}.
  The separate diagnostic script checks symbolic elimination and
  plots locally computed stationary values.
\item \textbf{Exact exponent comparisons.}
  \path{code/verify_order_constants.py} uses rational atanh-series
  bounds for logarithms to check the strict comparisons between the
  exponents in Section~\ref{ord:sec:main}. These finite sign checks
  support the stated separation of scales; the uniform analytic
  arguments remain part of the proof.
\item \textbf{Appell and transfer arithmetic.}
  \path{code/verify_mesh.py} records exact algebraic and rational
  checks used in the mesh constants and computes the sufficient
  integer thresholds without rounding them downward. Its finite
  examples are identified as checks of formulas, not a replacement
  for the argument valid at every degree.
\item \textbf{The frozen weight-fifteen certificate.}
  \path{code/verify_s14.py} starts from the integer vector and uses
  Python integers and rational fractions throughout. Its finite
  Euler sums, elementary constant intervals, and analytic tails
  produce the shipped interval for $R_{14}$ and verify
  $|R_{14}|<10^{-775}$. Both endpoints and all basket intervals
  are replayed exactly. The certificate contains zero and does not
  prove $R_{14}=0$.
\item \textbf{Independent diagnostics.}
  Symbolic differentiation verifies the negative-order kernels,
  and separate Dirichlet and Mellin evaluations check their displayed
  decimals. The mixed-generator program checks rational samples,
  integer samples, pole finite parts, and resummation tails.
  A separate Mellin calculation checks the frozen $S_{14}$ vector
  without using its discovery summation. These high-precision
  computations are corroboration, with their precision and residuals
  recorded in the data files.
\end{enumerate}

The default replay uses only the Python standard library for its exact
certificates. Optional numerical and symbolic diagnostics use
\texttt{mpmath}, \texttt{sympy}, and \texttt{matplotlib}; the versions
used in this delivery are recorded. The main \path{article.tex}
builds with pdfLaTeX and standard mathematics packages.
The source revision, incoming archive hashes, relevant source hashes,
verification commands, and PDF review are included in the package.
No repository modification or publication is part of this delivery.

\section{Further research questions}
\label{sec:research}

The results above remove several qualitative obstacles, but they also
identify where the present methods lose information. The following
questions are ordered by their relation to the proved theorems and by
the kind of additional argument they require.

\subsection{Finite-index extremizers and the order boundary}

\paragraph{1. Is every $K_N$ uniquely maximized, and is $M_N$ strictly
decreasing for every $N\ge2$?}
The exact $N=2$ elimination and the analytic large-$N$ theorem leave
a finite but presently unbounded intermediate region. One route is
to obtain a quantitative neighborhood and starting index for the
implicit-function argument, then certify the remaining rational
kernels by elimination or interval subdivision. A more satisfying
answer would be a direct comparison principle in $N$.
Pointwise monotonicity of a fixed-order normalized remainder is
unavailable in general, and monotonicity of the auxiliary envelope
does not alone compare its underlying maxima.

\paragraph{2. Does exact axis reduction hold at every truncation?}
For sufficiently large $N$, every positive outer order is strictly
inferior to the axis. The full finite-index question is to decide
whether $C_N=\max_b A_N(b)$ holds for all $N$, and whether the
axis maximum is always unique. If it fails, identify the first
interior maximizer and its transition. The uniform comparison
\eqref{ord:eq:axis-comparison} is a useful starting inequality, but
its positive and negative terms are too coarse at modest indices
to decide the question without further work.

\paragraph{3. Make the eventual statements effective.}
Produce explicit starting indices for axis reduction, axis uniqueness,
kernel uniqueness, and the respective monotonicity results.
The present proofs give concrete inequalities but leave several
\emph{sufficiently large} transitions unevaluated. Quantitative gamma
tail estimates and explicit bounds on derivatives of the limiting
kernels could turn them into finite gates. The practical challenge is
to obtain useful indices rather than enormous formal upper bounds.

\paragraph{4. Describe the threshold regimes beyond the maximizing
order.}
The Mellin integral in Lemma~\ref{ord:lem:threshold} converges for
$3/2<\alpha<5/2$. At the endpoints, a Taylor subtraction collides
with the integrability boundary. Determine uniform formulas as
$b/\log N$ approaches these values at a rate depending on $N$,
and determine the subsequent subtraction hierarchy. The poles of
the continued gamma factor suggest where logarithmic transition
terms should enter, but the integral must be treated before any
such expansion is asserted.

\paragraph{5. Obtain usable remainder bounds at the second scale.}
The inverse-log coefficient near the maximizing slope is large,
so a few formal terms are not a reliable moderate-$N$ approximation.
Find an expression retaining the threshold integral, or another
uniform approximation, with a computable relative error.
This would connect the exact extremizer theorem with practical
Euler truncation budgets. It would also allow meaningful numerical
tests of the optimizer displacement well before the extremely
late asymptotic regime.

\subsection{Appell spacing and joint zero geometry}

\paragraph{6. Determine the true minimum-gap asymptotic.}
The bound $1/(10n^2)$ answers the polynomial-separation question,
but it is not expected to describe every feature of the root
configuration. Determine the actual order and leading constant of
$\mesh(Q_n)$, and locate the pair of roots where the minimum occurs.
Bulk zero distributions alone need not control the smallest spacing,
especially if it occurs at an edge. A result resolving this distinction
would feed directly into the scale-free transfer theorem.

\paragraph{7. Improve the remote-block method.}
The proof uses $3n$ factors near $15n^2$ because that choice allows
elementary Laguerre estimates with transparent constants. Optimize
the block length, its position, and the ordered-root comparison.
Can information on correlated root motion reduce the loss from
the worst displacement of one endpoint? Can several blocks yield
a stronger power of $n$? The mesh-preservation theorem should
remain separate from any unsupported monotonicity in an individual
factor coefficient.

\paragraph{8. Find the optimal saturation scale and first offsets
when $n$ grows.}
The present sufficient threshold is $O(n^4\log^2 n)$, uniformly
over $\rho$ and every root. It is not a lower bound on the first
derivative order where saturation actually occurs. Determine the
true scale, its dependence on the branch, and whether the harmonic
factor is essential. Then strengthen the leading relative location
to a uniform additive offset and a controlled higher expansion.
This requires bounds on root-conditioned derivatives, not merely
separation and total zero count.

\paragraph{9. Extend the transfer theorem to other positive weights.}
Its essential input is a bound on the logarithmic derivative of the
positive weight, together with a real-rooted factor in $\log t$.
Identify other polylogarithmic and zeta derivative families with
that structure, including deformations whose logarithmic Lipschitz
constant grows with the degree. Determine when the resulting
constants still permit a polynomial joint-growth range.
Sign-changing weights require additional control because the
global variation bound may then change.

\subsection{Sign transitions and exact period reductions}

\paragraph{10. Classify every negative integer outer order.}
The finite product formula gives an explicit rational Gaussian
kernel for every $m$. Determine its number of roots in $(0,1)$,
then determine which permitted sign changes survive its gamma
average as $b$ varies. The $m=5$ example already separates the
one-crossing mechanism from the general problem.
Possible tools include reciprocal-polynomial substitutions,
Eulerian polynomial recurrences, and Mellin variation arguments.
A root count for the rational kernel is only an upper bound for
the number of transitions in the averaged function.

\paragraph{11. Construct a generating identity for the fixed-weight
mixed Gaussian reductions.}
The exact identities for $\mathcal S$ and $\mathcal O$ evaluate
infinite sums across weights. They do not supply the finite
weight-by-weight relation needed for $S_6$ or $S_{14}$.
The central problem is to find a bivariate generating relation,
a differential equation with enough boundary data, or an explicit
iterated-integral reduction whose coefficients are the desired
finite identities. Rational sampling may constrain such a
construction and indicate which cyclotomic levels it naturally uses.

\paragraph{12. Prove or refute the frozen $S_{14}$ relation.}
The interval in Proposition~\ref{mix:prop:proximity} is a precise
target for symbolic work. A proof should exhibit an exact algebraic
or iterated-integral relation, including endpoint terms and branch
conventions. Further digits would only strengthen proximity.
If the current finite relation space is insufficient, specify the
enlarged alphabet or class of relations and test the frozen vector
there. Do not infer period independence from a restricted rank
calculation.

\paragraph{13. Determine the arithmetic structure of the coefficient
vectors.}
The previously recorded $S_{10},S_{12}$ candidates and the present
$S_{14}$ vector have a common basket pattern. Establish whether
there is a uniform rational formula, recurrence, or generating
function for these coefficients, and explain the persistent
$-2\beta(2m)\log2$ term where it occurs. Searches should freeze
their basket and vector before independent evaluation and should
report exact residual intervals with an explicit normalization.
Any conjectured coefficient law needs testing outside the weights
used to infer it.

\subsection{Reproducible and formal continuations}

\paragraph{14. Formalize small reusable proof components.}
Several arguments here are natural proof-assistant targets:
the divided-difference expectation, the elementary global
maximum criterion, the rational $N=2$ elimination certificate,
the ordered-root comparison for $1+c\partial$, and the gamma
log-moment bound. They can be separated from the heavier complex
analytic continuation machinery. A formal integration plan should
distinguish algebraic certificate checking from formalized analysis
and should retain the exact source revision for inherited
dependencies.

The first mathematical priorities are an effective finite-index
extremizer classification, a sharper Appell mesh theorem, and an
exact relation producing the fixed-weight mixed Gaussian identities.
Each has a concrete obstruction exposed by the present proofs,
and none is resolved by additional numerical precision alone.
